% !TeX root = ../高等代数习题（谢启鸿）笔记.tex
% 由主文档通过 \input 载入；请编译主文档。
\setcounter{chapter}{2} % 保留原书章号。
\chapter{线性空间与线性方程组}
\label{chap:linear-spaces-systems}

% 研读来源：高等代数/textbooks/高等代数习题（谢启鸿）/高代习题-3.pdf，
% 第3章，书页116--191；按原书十节组织，整合基本结论、典型方法与训练题。
线性方程组的消元过程不仅给出未知数的取值，还揭示了向量之间有多少独立关系、一个空间需要多少个自由参数以及矩阵在乘法中损失了多少维信息。本章把这些问题统一在线性空间的语言中：先由线性表示建立基与维数，再用坐标把抽象问题化成矩阵计算，最后以秩和解空间解释方程组的可解性及其应用。学习时应把握两条相互联系的路线：具体计算依靠消元，结构证明依靠生成空间、线性无关性及维数。

除特别说明外，$\mathbb F$ 表示实数域或复数域的子域，各线性空间均在同一数域 $\mathbb F$ 上，讨论基、维数和维数公式时默认空间有限维。坐标向量统一写成列向量，$M_{m\times n}(\mathbb F)=M_{m,n}(\mathbb F)$ 表示 $m\times n$ 矩阵全体，$M_n(\mathbb F)=M_{n\times n}(\mathbb F)$，$I_n$ 表示 $n$ 阶单位矩阵。$A^{\mathsf T}$ 表示转置，$A^{\mathsf H}=\overline A^{\mathsf T}$ 表示共轭转置，$\operatorname{adj}(A)$ 表示伴随矩阵；三种记号所表示的运算不同。

\section{基本概念}
\label{sec:laex-basic}

\subsection{线性空间、线性表示与线性无关}

\begin{definition}
\label{def:laex-vector-space}
设 $V$ 是非空集合，在 $V$ 上定义了加法，在 $\mathbb F\times V$ 上定义了数乘。若任意 $u,v,w\in V$ 与 $a,b\in\mathbb F$ 满足
\[
\begin{gathered}
u+v=v+u,\qquad (u+v)+w=u+(v+w),\\
\text{存在 }0\in V\text{ 使 }v+0=v,\qquad
\text{对每个 }v\text{ 存在 }-v\text{ 使 }v+(-v)=0,\\
1v=v,\qquad a(bv)=(ab)v,\\
a(u+v)=au+av,\qquad (a+b)v=av+bv,
\end{gathered}
\]
则称 $V$ 为 $\mathbb F$ 上的\emph{线性空间（vector space）}，其元素称为向量。
\end{definition}

线性空间由集合、加法、数乘和标量域共同确定，不能只看元素的外观。通常运算下的 $\mathbb F^n$、矩阵空间、多项式空间以及实连续函数空间都是典型例子；在这些例子中，向量可以是有序数组、矩阵、多项式或函数。次数恰为 $n$ 的多项式全体则不是线性空间，因为它不含零多项式，而且两个这样的多项式相加可能降低次数。

\begin{definition}
\label{def:laex-span-independence}
给定向量 $v_1,\ldots,v_m\in V$，形如 $a_1v_1+\cdots+a_mv_m$ 的向量称为它们的\emph{线性组合（linear combination）}，其中 $a_i\in\mathbb F$。所有这些组合组成的集合记为
\[
\operatorname{span}(v_1,\ldots,v_m)
=\left\{\sum_{i=1}^m a_iv_i:a_i\in\mathbb F\right\},
\]
称为该向量组的生成空间。若等式 $a_1v_1+\cdots+a_mv_m=0$ 只能由 $a_1=\cdots=a_m=0$ 得到，则称向量组\emph{线性无关（linearly independent）}；否则称其线性相关。约定空向量组线性无关，且其生成空间为 $\{0\}$。
\end{definition}

这里的等式是空间中的向量等式。例如，函数的线性关系必须对定义域中的每一点成立，多项式的线性关系必须是恒等式，不能只在未经说明的若干个点上验证。一个非空向量组只要含零向量就线性相关，但不含零向量远不能保证线性无关。

\begin{proposition}
\label{prop:laex-dependence-basic}
有限向量组线性相关，当且仅当其中至少一个向量可以由其余向量线性表示。若 $v_1,\ldots,v_m$ 线性无关，则对任意 $v\in V$，加入 $v$ 后的向量组线性无关，当且仅当 $v\notin\operatorname{span}(v_1,\ldots,v_m)$。一个向量在给定向量组中的线性表示一旦存在，其表示系数唯一，当且仅当该向量组线性无关。
\end{proposition}

\begin{proof}
非平凡关系 $\sum_i a_iv_i=0$ 中至少有一个 $a_j\ne0$，移项后除以 $a_j$ 就得到 $v_j$ 的表示；反过来，将这样的表示移到等号一侧就得到非平凡关系。若 $\sum_i a_iv_i+bv=0$，当 $b=0$ 时，原向量组无关性迫使所有 $a_i=0$；因此加入新向量后出现非平凡关系，当且仅当可以从 $b\ne0$ 的关系中解出 $v$。最后，两种表示相减得到齐次关系；若存在非平凡齐次关系，将其加在已有表示上就会产生另一组系数。
\end{proof}

\subsection{基、维数、子空间与矩阵的秩}

\begin{definition}
\label{def:laex-basis-rank}
一个向量组的\emph{极大线性无关组（maximal linearly independent subset）}是其中既线性无关、又能生成原向量组中全部向量的部分组。其所含向量个数称为原向量组的\emph{秩（rank）}。若线性无关组能生成整个空间 $V$，则称其为 $V$ 的一组\emph{基（basis）}；有限基所含向量个数称为 $V$ 的\emph{维数（dimension）}，记为 $\dim_{\mathbb F}V$，标量域明确时简记为 $\dim V$。零空间的基为空组，维数为零。
\end{definition}

这些定义中的个数不依赖所选的极大无关组，依据是如下基本比较原理：若 $r$ 个线性无关向量均可由 $s$ 个向量生成，则 $r\le s$。事实上，把这 $r$ 个向量的表示系数排成 $s\times r$ 矩阵；若 $r>s$，对相应的齐次方程组消元，主元至多有 $s$ 个，必有自由未知数，因而存在非零系数使这 $r$ 个向量的组合为零，产生矛盾。将此原理分别用于两组基，就得到它们具有相同个数。

\begin{definition}
\label{def:laex-subspace}
若 $U\subseteq V$ 在原有加法和数乘下也是线性空间，则称 $U$ 为 $V$ 的\emph{子空间（subspace）}。判定时只需检查 $U\ne\varnothing$，且对任意 $u,v\in U$ 与 $a,b\in\mathbb F$ 有 $au+bv\in U$。两个子空间的交与和分别为
\[
U\cap W=\{v:v\in U,\ v\in W\},\qquad
U+W=\{u+w:u\in U,\ w\in W\}.
\]
两者都是子空间。若 $U\cap W=\{0\}$，则称这个和为\emph{直和（direct sum）}，记为 $U\oplus W$。
\end{definition}

矩阵 $A\in M_{m\times n}(\mathbb F)$ 的列空间是其列向量的生成空间，记为 $\operatorname{im}A=\{Ax:x\in\mathbb F^n\}$；核空间记为 $\ker A=\{x\in\mathbb F^n:Ax=0\}$。矩阵的秩定义为列空间的维数，记为 $\operatorname{rank}A$。下一节将说明它也等于行向量组的秩，以及行阶梯形矩阵中主元的个数。由消元可知，若 $\operatorname{rank}A=r$，齐次方程组就有 $n-r$ 个自由未知数，依次令各自由未知数取标准单位向量得到其解空间的一组基，因此
\begin{equation}
\label{eq:laex-rank-nullity}
\dim\ker A+\operatorname{rank}A=n.
\end{equation}
这个等式称为\emph{秩与零度公式（rank--nullity formula）}；$n$ 是未知数个数，不是方程个数。

\section{向量的线性关系}
\label{sec:laex-linear-relations}

\subsection{消元、秩与极大无关组}

具体向量的线性关系可以转化为矩阵问题。将 $v_1,\ldots,v_s\in\mathbb F^n$ 按列组成 $A=(v_1,\ldots,v_s)$，则 $\sum_j c_jv_j=0$ 等价于 $Ac=0$。若向量原来写成行向量，仍可先转置再按列排列，以便统一使用下面的方法。

\begin{proposition}
\label{prop:laex-pivot-columns}
设 $R=PA$ 是 $A$ 经初等行变换得到的行阶梯形矩阵，其中 $P$ 可逆。则 $A$ 与 $R$ 的列向量之间具有完全相同的线性关系。若 $R$ 有 $r$ 个非零行，主元列的列指标为 $j_1,\ldots,j_r$，则 $\operatorname{rank}A=r$，原矩阵中第 $j_1,\ldots,j_r$ 列组成列向量组的一个极大无关组。矩阵的行秩与列秩相等。
\end{proposition}

\begin{proof}
由 $P$ 可逆，$Ac=0$ 当且仅当 $Rc=0$。$R$ 的主元列在前 $r$ 行构成一个对角元非零的上三角矩阵，故这 $r$ 列无关；而所有列的后续分量均为零，其生成空间维数至多为 $r$，所以主元列是极大无关组。$R$ 的 $r$ 个非零行也线性无关：依次考察各主元位置，任何为零的行向量线性组合都迫使系数逐个为零。初等行变换保持行空间不变，故 $A$ 的行秩也为 $r$。
\end{proof}

\begin{notealert}[type=caution,title={返回原矩阵取列向量}]
求原列向量组的极大无关组时，应使用行阶梯形矩阵确定列指标，再返回原矩阵取相应列。行变换保持列之间的关系，却通常改变各列本身；任意列变换只保证矩阵的秩不变，不能保证原列指标仍有上述意义。
\end{notealert}

\begin{example}
\label{ex:laex-maximal-subset}
求向量组
\[
\begin{aligned}
v_1&=(1,1,2,2,1)^{\mathsf T},&v_2&=(0,2,1,5,-1)^{\mathsf T},\\
v_3&=(2,0,3,-1,3)^{\mathsf T},&v_4&=(1,1,0,4,-1)^{\mathsf T}
\end{aligned}
\]
的秩与一个极大无关组。
\end{example}

\begin{solve}
先观察到 $v_3=2v_1-v_2$，所以原向量组可由 $v_1,v_2,v_4$ 生成。取这三列的前三个分量得到子矩阵
\[
\begin{pmatrix}1&0&1\\1&2&1\\2&1&0\end{pmatrix},
\]
其行列式为 $-4\ne0$，故 $v_1,v_2,v_4$ 无关。因此秩为 $3$，它们是一个极大无关组。又因为 $v_3=2v_1-v_2$ 中两个系数均非零，也可以用 $v_3$ 替换 $v_1$ 或 $v_2$，得到其他极大无关组；向量的选择可以不同，个数则始终为 $3$。
\end{solve}

判定 $b$ 是否可由 $v_1,\ldots,v_s$ 线性表示，就是判断 $Ax=b$ 是否有解，等价条件为
\begin{equation}
\label{eq:laex-linear-representation-rank}
\operatorname{rank}(A\mid b)=\operatorname{rank}A.
\end{equation}
若还需要具体系数，就应对增广矩阵消元，而不能只比较两个秩。

\begin{example}
\label{ex:laex-representation}
设 $v_1=(1,0,-1,3)^{\mathsf T}$，$v_2=(2,1,0,1)^{\mathsf T}$，$v_3=(0,-2,1,1)^{\mathsf T}$，判断 $b=(5,4,-2,4)^{\mathsf T}$ 能否由它们表示，并说明表示是否唯一。
\end{example}

\begin{solve}
直接代入可得 $b=v_1+2v_2-v_3$，故表示存在。前三个分量组成的系数矩阵行列式为 $5\ne0$，所以三列线性无关，由命题~\ref{prop:laex-dependence-basic} 可知这组表示系数唯一。
\end{solve}

\subsection{由定义证明线性关系}

抽象向量没有现成分量时，通常先设一个线性组合为零，再利用已知的线性无关性逐步消去系数。一个常用推论是：若 $b$ 能由 $v_1,\ldots,v_m$ 表示，却不能由其中任何更少的向量表示，则这 $m$ 个向量无关。否则其中一个向量可由其余向量表示，将它代回 $b$ 的表达式就减少了所需向量个数。

\begin{proposition}
\label{prop:laex-generating-rank}
若向量组 $B$ 中每个向量均可由向量组 $A$ 线性表示，则
\[
\operatorname{span}B\subseteq\operatorname{span}A,
\qquad \operatorname{rank}B\le\operatorname{rank}A.
\]
两组向量互相线性表示，称为\emph{等价向量组}；这等价于它们生成同一个子空间。若已知其中一个生成空间包含于另一个，则两组秩相等也足以推出等价。
\end{proposition}

\begin{proof}
把 $B$ 中各向量的表达式代入任意线性组合，就得到生成空间的包含。取 $B$ 的一个极大无关组，用 $A$ 的极大无关组表示它，由基本比较原理得到秩的不等式。最后，两个有限维子空间若有包含关系且维数相同，就必相等：否则可从较大空间取一个不在较小空间中的向量，加入后增加一维，产生矛盾。
\end{proof}

\begin{example}
\label{ex:laex-coefficient-matrix}
已知 $v_1,\ldots,v_m$ 线性无关，令
\[
w_j=\sum_{i=1}^m c_{ij}v_i\quad(1\le j\le s),
\qquad C=(c_{ij})\in M_{m\times s}(\mathbb F).
\]
证明 $\operatorname{rank}(w_1,\ldots,w_s)=\operatorname{rank}C$。
\end{example}

\begin{solve}
对任意 $t\in\mathbb F^s$，有
\[
\sum_{j=1}^s t_jw_j
=\sum_{i=1}^m\left(\sum_{j=1}^s c_{ij}t_j\right)v_i=0
\quad\Longleftrightarrow\quad Ct=0.
\]
因此两个向量组有完全相同的系数关系，极大无关组的指标也相同，秩自然相等。特别地，当 $s=m$ 时，新向量组无关当且仅当 $\det C\ne0$。若不假设原向量组无关，一般只能断言新向量组的秩不超过 $\operatorname{rank}C$，不能直接取等号。
\end{solve}

\begin{example}
\label{ex:laex-cyclic-vectors}
设 $v_1,\ldots,v_s$ 是一组基，$s>1$，并令 $w_i=v_i+v_{i+1}$，其中 $v_{s+1}=v_1$。讨论 $w_1,\ldots,w_s$ 的线性关系。
\end{example}

\begin{solve}
关系 $\sum_i c_iw_i=0$ 在原基下等价于 $c_i+c_{i-1}=0$，下标按模 $s$ 解释。因此系数沿环交替变号。若 $s$ 为奇数，绕一周得到 $c_1=-c_1$，因 $\mathbb F$ 的特征为零，只能所有系数为零，新向量组是一组基；若 $s$ 为偶数，则关系空间恰为一维，关系可取 $w_1-w_2+\cdots-w_s=0$，秩为 $s-1$。
\end{solve}

\begin{notealert}[type=caution,title={区分不同层次的线性关系}]
两两无关不能推出整个向量组无关，例如 $e_1,e_2,e_1+e_2$。两组向量的秩相同也不能推出它们等价，例如 $\{e_1\}$ 与 $\{e_2\}$。矩阵相抵只要求同型矩阵的秩相同，而向量组等价要求生成空间相同，不能混用这两种“等价”。
\end{notealert}

\section{线性空间及其基}
\label{sec:laex-spaces-bases}

\subsection{基的判定与扩张}

\begin{theorem}
\label{thm:laex-basis-extension}
设 $\dim V=n$。$V$ 中任意线性无关组都可以扩张为一组基；任意有限生成组都可以删减为一组基。因此，恰含 $n$ 个向量的向量组只要满足“线性无关”或“生成 $V$”中的任一条件，就是 $V$ 的一组基。
\end{theorem}

\begin{proof}
从无关组出发，若尚未生成 $V$，就在其生成空间外选取一个向量，由命题~\ref{prop:laex-dependence-basic} 得到更大的无关组。基本比较原理保证无关向量个数不超过 $n$，故有限步后得到基。从有限生成组出发，若其相关，就删除一个可由其余向量表示的向量，生成空间不变，直至得到无关组。若原向量组已含 $n$ 个向量，无关组不能继续扩张，而生成组也不能删去任何向量，否则基的个数会小于 $n$。
\end{proof}

\begin{example}
\label{ex:laex-polynomial-bases}
记 $\mathcal P_n(\mathbb F)=\{f\in\mathbb F[x]:\deg f\le n\}\cup\{0\}$，其中 $n\ge0$。说明它的维数，并在 $n+1$ 个互异节点 $a_0,\ldots,a_n\in\mathbb F$ 上构造一组便于求坐标的基。
\end{example}

\begin{solve}
按系数比较可知 $1,x,\ldots,x^n$ 既无关又生成该空间，所以维数为 $n+1$。定义
\[
\ell_j(x)=\prod_{\substack{0\le k\le n\\k\ne j}}
\frac{x-a_k}{a_j-a_k}\quad(0\le j\le n),
\]
则 $\ell_j(a_i)=\delta_{ij}$。若 $\sum_j c_j\ell_j=0$，依次代入各节点得到所有 $c_j=0$，故这 $n+1$ 个多项式是一组基。任意 $f\in\mathcal P_n(\mathbb F)$ 因而满足
\[
f(x)=\sum_{j=0}^n f(a_j)\ell_j(x).
\]
这就是拉格朗日插值公式（Lagrange interpolation formula），其坐标是节点处的函数值。次数不超过 $n$ 的限制保证这些节点值足以唯一确定多项式。
\end{solve}

\begin{example}
\label{ex:laex-matrix-bases}
求全体 $m\times n$ 矩阵、全体 $n$ 阶上三角矩阵、对称矩阵和反对称矩阵空间的维数。
\end{example}

\begin{solve}
令 $E_{ij}$ 是仅在第 $(i,j)$ 个位置为 $1$、其余位置均为零的矩阵。全体 $E_{ij}$ 构成 $M_{m\times n}(\mathbb F)$ 的基，维数为 $mn$。上三角矩阵空间的基为 $\{E_{ij}:i\le j\}$，维数为 $n(n+1)/2$；对称矩阵空间的基为
\[
\{E_{ii}:1\le i\le n\}\cup
\{E_{ij}+E_{ji}:1\le i<j\le n\},
\]
维数也为 $n(n+1)/2$；反对称矩阵空间的基为 $\{E_{ij}-E_{ji}:i<j\}$，维数为 $n(n-1)/2$。这些组的生成性由逐项展开得到，无关性则由比较各独立位置的元素得到。
\end{solve}

\subsection{标量域与无限维现象}

\begin{example}
\label{ex:laex-hermitian}
令 $\mathcal H_n=\{A\in M_n(\mathbb C):A^{\mathsf H}=A\}$，$\mathcal K_n=\{A\in M_n(\mathbb C):A^{\mathsf H}=-A\}$，分别为埃尔米特（Hermitian）矩阵与斜埃尔米特（skew-Hermitian）矩阵全体。说明它们作为实线性空间的维数。
\end{example}

\begin{solve}
这两个集合对加法和实数数乘封闭。$\mathcal H_n$ 的一组实基为
\[
\{E_{ii}\}\cup\{E_{ij}+E_{ji}:i<j\}
\cup\{\mathrm i(E_{ij}-E_{ji}):i<j\},
\]
所以 $\dim_{\mathbb R}\mathcal H_n=n+2\binom n2=n^2$。映射 $A\mapsto\mathrm iA$ 将此基送到 $\mathcal K_n$ 的一组实基，故后者维数也为 $n^2$。两者通常都不是复线性空间：非零 Hermite 矩阵乘 $\mathrm i$ 后成为斜 Hermite 矩阵，已不在原集合中。
\end{solve}

同一个集合在不同标量域上的维数可能不同。例如 $\mathbb C$ 在自身上以 $1$ 为基，维数为 $1$；在 $\mathbb R$ 上则以 $1,\mathrm i$ 为基，维数为 $2$。遇到数域扩张时，必须把允许使用的系数写清楚。

\begin{example}
\label{ex:laex-cubic-field}
设 $\theta=\sqrt[3]{2}$，证明 $W=\{a+b\theta+c\theta^2:a,b,c\in\mathbb Q\}$ 是数域，并求其在 $\mathbb Q$ 上的维数。
\end{example}

\begin{solve}
多项式 $x^3-2$ 在 $\mathbb Q$ 上无有理根，因其次数为三，所以不可约。若一个次数至多为二的非零有理系数多项式也以 $\theta$ 为根，它与 $x^3-2$ 互素，由多项式的 Bézout 等式就会推出 $1=0$，故不可能。因此 $1,\theta,\theta^2$ 无关，$\dim_{\mathbb Q}W=3$。由 $\theta^3=2$ 可知 $W$ 对乘法封闭。对任意 $0\ne u\in W$，映射 $v\mapsto uv$ 是 $W$ 到自身的单射线性映射，因为 $W\subset\mathbb R$ 中无零因子；有限维且维数相同使它也是满射。于是存在 $v\in W$ 使 $uv=1$，证明非零元的逆仍在 $W$ 中，故 $W$ 为数域。
\end{solve}

\begin{proposition}
\label{prop:laex-scalar-restriction}
设 $\mathbb F\subseteq\mathbb K$，$\mathbb K$ 在 $\mathbb F$ 上有基 $a_1,\ldots,a_m$，而 $V$ 在 $\mathbb K$ 上有基 $e_1,\ldots,e_n$。则 $\{a_ie_j:1\le i\le m,\ 1\le j\le n\}$ 是 $V$ 在 $\mathbb F$ 上的一组基，特别地
\[
\dim_{\mathbb F}V=m\dim_{\mathbb K}V=mn.
\]
\end{proposition}

\begin{proof}
把任意 $v=\sum_j b_je_j$ 中的每个 $b_j\in\mathbb K$ 展开为 $\sum_i c_{ij}a_i$，就得到这些乘积向量的生成性。若 $\sum_{i,j}c_{ij}a_ie_j=0$，先用 $e_j$ 的 $\mathbb K$ 线性无关性得到 $\sum_i c_{ij}a_i=0$，再用 $a_i$ 的 $\mathbb F$ 线性无关性得到所有 $c_{ij}=0$。
\end{proof}

无限维不能仅理解为“元素有无穷多个”，因为任何非零有限维数域上线性空间都已经有无穷多个元素。要证明空间不是有限维，应构造任意多个线性无关向量。例如 $C(\mathbb R)$ 包含任意长的无关组 $1,x,\ldots,x^N$；趋于零的实数列空间 $c_0$ 包含标准单位数列 $e_1,\ldots,e_N$，它们也无关。无限向量族的线性无关是指其每个有限子族都无关，代数意义的基只允许\emph[caution]{有限线性组合}；因此不能把任意无限级数自动当作基展开。后文使用的维数相减与有限步扩张论证均限于已声明的有限维情形。

\section{线性同构和几何问题代数化}
\label{sec:laex-isomorphism}

\begin{definition}
\label{def:laex-isomorphism}
设 $V,W$ 是同一数域 $\mathbb F$ 上的线性空间。若映射 $T:V\to W$ 为双射，且对任意 $u,v\in V$、$a,b\in\mathbb F$ 满足 $T(au+bv)=aT(u)+bT(v)$，则称 $T$ 为\emph{线性同构（linear isomorphism）}，记为 $V\cong W$。
\end{definition}

同构保持全部有限线性关系，因为 $\sum_i c_iT(v_i)=T(\sum_i c_iv_i)$，再由单射性可知左端为零当且仅当原组合为零。因此同构保持线性无关性、生成关系、向量组的秩及空间维数。同一数域上两个有限维线性空间同构，当且仅当它们维数相同：必要性由同构送基到基得到；充分性则由把一组基逐个对应到另一组基，并按线性组合延拓得到。

\begin{theorem}
\label{thm:laex-coordinate-isomorphism}
给定有序基 $E=(e_1,\ldots,e_n)$，对 $v=\sum_i x_ie_i$ 定义坐标列向量 $[v]_E=(x_1,\ldots,x_n)^{\mathsf T}$，则映射 $v\mapsto[v]_E$ 是 $V$ 到 $\mathbb F^n$ 的线性同构。因此，抽象向量组的秩等于其坐标列组成的矩阵的秩，且极大无关组的列指标相同。
\end{theorem}

\begin{proof}
基展开的存在性与唯一性分别保证映射为双射，两个基展开相加或同乘一个标量则保证线性。其余结论由同构保持线性关系得到。
\end{proof}

\emph[tip]{几何问题代数化}的关键是先固定同一组基，把每个向量转换成坐标列，再用矩阵运算求解，最后将答案解释回原空间。不同向量若使用不同基，所得坐标列就不能直接拼成同一个矩阵。例如函数或多项式的坐标既可以是通常系数，也可以是在适当基下的节点值；选择坐标的目的，是使原问题的线性关系更容易计算。

\begin{example}
\label{ex:laex-vandermonde-basis}
设 $E=(e_1,\ldots,e_n)$ 是 $V$ 的一组基，$a_1,\ldots,a_n\in\mathbb F$ 互不相同，令
\[
v_j=e_1+a_je_2+\cdots+a_j^{n-1}e_n\quad(1\le j\le n).
\]
证明 $v_1,\ldots,v_n$ 也是一组基。
\end{example}

\begin{solve}
坐标矩阵是范德蒙德矩阵（Vandermonde matrix），其行列式为
\[
\det\begin{pmatrix}
1&\cdots&1\\a_1&\cdots&a_n\\\vdots&&\vdots\\a_1^{n-1}&\cdots&a_n^{n-1}
\end{pmatrix}
=\prod_{1\le i<j\le n}(a_j-a_i)\ne0.
\]
因此坐标列无关，原向量也无关，再由定理~\ref{thm:laex-basis-extension} 得到结论。这种构造把任意 $n$ 个互异标量变成一组基，后面可用于避开有限多个真子空间。
\end{solve}

\begin{example}
\label{ex:laex-unusual-operations}
在 $\mathbb R_{>0}$ 上定义 $u\oplus v=uv$、$a\odot u=u^a$。另在 $\mathbb R^2$ 上定义
\[
(u,v)\oplus(s,t)=(u+s,v+t+us),\qquad
a\odot(u,v)=\left(au,av+\frac{a(a-1)}2u^2\right).
\]
说明这两组运算都使相应集合成为实线性空间。
\end{example}

\begin{solve}
第一种情形的映射 $T:\mathbb R\to\mathbb R_{>0}$，$T(x)=e^x$ 为双射，满足 $T(x+y)=T(x)\oplus T(y)$、$T(ax)=a\odot T(x)$。因此新运算由通常实线性空间的运算传递而来，八条公理均成立；其零向量是数 $1$，加法逆元是通常乘法的倒数。

第二种情形使用双射 $S(x,y)=(x,y+x^2/2)$，其逆为 $S^{-1}(u,v)=(u,v-u^2/2)$。展开平方可验证 $S(z_1+z_2)=S(z_1)\oplus S(z_2)$ 与 $S(az)=a\odot S(z)$，因而同样构成线性空间。所谓“线性”是相对于所定义的加法和数乘而言，不能仅凭公式中出现平方或乘方就否定线性空间结构。
\end{solve}

\section{基变换与过渡矩阵}
\label{sec:laex-change-basis}

\subsection{列的含义与坐标变换方向}

\begin{definition}
\label{def:laex-transition}
设 $E=(e_1,\ldots,e_n)$ 与 $F=(f_1,\ldots,f_n)$ 是 $V$ 的两组有序基，令
\[
f_j=\sum_{i=1}^n p_{ij}e_i,\qquad P=(p_{ij}).
\]
称 $P$ 为从基 $E$ 到基 $F$ 的\emph{过渡矩阵（transition matrix）}。其第 $j$ 列就是 $[f_j]_E$，形式上记为 $F=EP$。
\end{definition}

若 $x=[v]_E$、$y=[v]_F$，则 $v=Ex=Fy=EPy$，由基展开唯一性得到
\begin{equation}
\label{eq:laex-coordinate-change}
x=Py,\qquad y=P^{-1}x.
\end{equation}
因此“从旧基到新基的过渡矩阵”将新坐标变成旧坐标；要把旧坐标变成新坐标，应使用其逆矩阵。$P$ 可逆，因为它的列是另一组基在 $E$ 下的坐标。若第三组基满足 $G=FQ$，则 $G=E(PQ)$，故过渡矩阵沿基变换次序相乘。

具体计算时，若基 $E,F$ 在某一共同参考基下的坐标矩阵分别为 $A,B$，则 $B=AP$，从而 $P=A^{-1}B$。可以对 $(A\mid B)$ 进行初等行变换，将左半块化为 $I_n$，右半块便同时化为 $P$，无须分别求出每个新基向量的系数。这里 $A,B$ 必须按列排列基向量。

\subsection{平移幂基与多项式的 Taylor 展开}

\begin{example}
\label{ex:laex-taylor-basis}
在 $\mathcal P_n(\mathbb F)$ 中，取 $E=(1,x,\ldots,x^n)$ 与 $F=(1,x-a,\ldots,(x-a)^n)$。求过渡矩阵，并解释后一组基下的坐标。
\end{example}

\begin{solve}
为简化公式，以 $0,1,\ldots,n$ 为行列指标。二项式定理给出
\[
p_{ij}=\begin{cases}\binom ji(-a)^{j-i},&i\le j,\\0,&i>j.\end{cases}
\]
矩阵 $P$ 为对角元全为 $1$ 的上三角矩阵，因而 $F$ 确为一组基。逆变换来自 $x=(x-a)+a$，故 $P^{-1}$ 的对应元素为 $\binom ji a^{j-i}$，$i\le j$。

设 $f(x)=\sum_{j=0}^n b_j(x-a)^j$。对等式求 $k$ 次导数，再令 $x=a$，只剩 $k!b_k$，得到
\begin{equation}
\label{eq:laex-polynomial-taylor}
f(x)=\sum_{k=0}^n\frac{f^{(k)}(a)}{k!}(x-a)^k.
\end{equation}
所以 Taylor 公式在这里就是一次基变换：新坐标为各阶导数值除以相应阶乘。这是有限次数多项式的恒等式，不涉及无穷级数、收敛性或余项估计。
\end{solve}

\section{子空间与商空间}
\label{sec:laex-subspace-quotient}

\subsection{交、和及其计算}

\begin{theorem}
\label{thm:laex-dimension-sum}
设 $U,W$ 是有限维线性空间 $V$ 的子空间，则
\begin{equation}
\label{eq:laex-dimension-sum}
\dim(U+W)=\dim U+\dim W-\dim(U\cap W).
\end{equation}
\end{theorem}

\begin{proof}
取 $U\cap W$ 的基 $e_1,\ldots,e_k$，分别扩张成 $U$ 的基 $e_1,\ldots,e_k,u_1,\ldots,u_p$ 和 $W$ 的基 $e_1,\ldots,e_k,w_1,\ldots,w_q$。合并且只保留一份公共基，所得向量组生成 $U+W$。若
\[
\sum_i a_ie_i+\sum_j b_ju_j+\sum_l c_lw_l=0,
\]
则 $\sum_l c_lw_l\in U\cap W$，于是 $W$ 中基展开唯一性迫使所有 $c_l=0$，再由 $U$ 的基无关性得到所有 $a_i,b_j=0$。所以合并组是一组基，维数为 $k+p+q$，恰等于右端。
\end{proof}

若 $U=\operatorname{im}A$、$W=\operatorname{im}B$，则 $U+W$ 的基可从 $(A\mid B)$ 的原列中选择。求交空间时解
\begin{equation}
\label{eq:laex-intersection-system}
As=Bt\quad\Longleftrightarrow\quad (A\mid-B)\binom st=0,
\end{equation}
再把解参数映射为公共向量 $As$。当 $A,B$ 的列分别为 $U,W$ 的基时，这个映射在上述系数解空间上是双射，因此基础解系直接对应交空间的基；若两侧只是含冗余的生成组，则不同系数可能给出同一向量，还应剔除映射后的零向量和相关向量。

\begin{example}
\label{ex:laex-intersection}
设
\[
\begin{aligned}
a_1&=(1,0,-1,0)^{\mathsf T},&a_2&=(0,1,2,1)^{\mathsf T},&a_3&=(2,1,0,1)^{\mathsf T},\\
b_1&=(-1,1,1,1)^{\mathsf T},&b_2&=(1,-1,-3,-1)^{\mathsf T},&b_3&=(-1,1,-1,1)^{\mathsf T}.
\end{aligned}
\]
令 $U=\operatorname{span}(a_1,a_2,a_3)$，$W=\operatorname{span}(b_1,b_2,b_3)$，求交与和的基。
\end{example}

\begin{solve}
因为 $a_3=2a_1+a_2$、$b_3=2b_1+b_2$，且 $a_1,a_2$ 和 $b_1,b_2$ 分别无关，故两空间维数都为 $2$。又 $b_2=a_1-a_2\in U$，而 $b_1\notin U$：若用前两个分量确定其系数为 $-1,1$，第三分量应为 $3$，实际却为 $1$。因此
\[
U+W=\operatorname{span}(a_1,a_2,b_1),\qquad
\dim(U\cap W)=2+2-3=1.
\]
所以 $a_1,a_2,b_1$ 是和空间的一组基，$b_2$ 是交空间的一组基。
\end{solve}

\subsection{直和、补空间与有限个真子空间}

\begin{proposition}
\label{prop:laex-direct-sum}
设 $V_1,\ldots,V_k$ 是有限维子空间，$S=V_1+\cdots+V_k$。下列条件等价：每个 $v\in S$ 的分解 $v=v_1+\cdots+v_k$ 唯一；$\sum_i v_i=0$ 只能由所有 $v_i=0$ 得到；对每个 $j\ge2$ 有 $V_j\cap(V_1+\cdots+V_{j-1})=\{0\}$；各空间的基拼接后是 $S$ 的一组基；$\dim S=\sum_i\dim V_i$。成立时记为 $S=\bigoplus_{i=1}^k V_i$。
\end{proposition}

\begin{proof}
两种分解相减就是零向量的分解，故前两个条件等价。在一个非平凡零分解中取最后一个非零分量，便得到相应的交中有非零向量，反向移项也成立。把每个分量按其空间的基展开，零分解的唯一性正好等价于拼接组无关；该组总能生成 $S$，所以又等价于其个数等于 $\dim S$。
\end{proof}

\begin{example}
\label{ex:laex-direct-decompositions}
全体方阵可以唯一写成对称矩阵与反对称矩阵之和；$\mathbb F^n$ 也可以唯一写成常数分量向量与分量和为零的向量之和。
\end{example}

\begin{solve}
对方阵有
\[
A=\frac{A+A^{\mathsf T}}2+\frac{A-A^{\mathsf T}}2.
\]
两项分别对称与反对称，且同时具有两种性质的矩阵只能为零，故为直和。对 $x=(x_1,\ldots,x_n)^{\mathsf T}$，令 $\bar x=(x_1+\cdots+x_n)/n$、$\mathbf1=(1,\ldots,1)^{\mathsf T}$，则
\[
x=\bar x\mathbf1+(x-\bar x\mathbf1).
\]
后一项的分量和为零，常数分量向量若分量和为零也只能是零。因此两空间构成直和，维数分别为 $1$ 和 $n-1$。这里除以 $2$ 和 $n$ 使用了数域的特征为零这一条件。
\end{solve}

\begin{notealert}[type=caution,title={多于两个子空间时不能只检验两两相交为零}]
在 $\mathbb R^2$ 中，$\operatorname{span}(e_1)$、$\operatorname{span}(e_2)$、$\operatorname{span}(e_1+e_2)$ 两两交为零，但三者之和不是直和，因为 $e_1+e_2-(e_1+e_2)=0$ 给出非平凡分解。应检验某一子空间与前面各空间之和的交，或检验维数是否相加。
\end{notealert}

任意子空间 $U\subseteq V$ 都有\emph{补空间（complementary subspace）} $W$，使 $V=U\oplus W$：取 $U$ 的基并扩张为 $V$ 的基，由新增基向量生成 $W$ 即可。补空间通常不唯一，例如 $\mathbb R^2$ 中 $\operatorname{span}(e_1)$ 的补空间可以是任意 $\operatorname{span}(e_2+ae_1)$，$a\in\mathbb R$。补空间与集合补集不同，二者交为 $\{0\}$，而且 $V=U+W$ 表示每个向量都能相加分解，并非 $V=U\cup W$。外直和 $U\times W$ 以逐分量运算定义，当 $U,W$ 已在某空间中构成内直和时，$(u,w)\mapsto u+w$ 就给出两种构造之间的同构。

\begin{proposition}
\label{prop:laex-finite-union}
设 $\dim V=n\ge1$，标量域 $\mathbb F$ 无限。有限个真子空间不能覆盖 $V$；更强地，存在 $V$ 的一组基，其中每个基向量都避开这些真子空间。
\end{proposition}

\begin{proof}
固定基 $e_1,\ldots,e_n$，对每个 $t\in\mathbb F$ 构造 $v(t)=e_1+te_2+\cdots+t^{n-1}e_n$。由例~\ref{ex:laex-vandermonde-basis}，任意 $n$ 个互异参数对应的向量组成一组基，故一个真子空间至多包含其中 $n-1$ 个不同参数对应的向量。有限个真子空间只排除了有限多个参数，从无限域中另选 $n$ 个互异参数就得到所需基。有限域上这一步失效，例如 $\mathbb F_2^2$ 是其三个一维子空间的并。
\end{proof}

\subsection{商空间及其维数}

\begin{definition}
\label{def:laex-quotient}
设 $U$ 是 $V$ 的子空间，$v+U=\{v+u:u\in U\}$ 称为 $U$ 的陪集（coset）。全体陪集构成的集合记为 $V/U$，并定义
\[
(v+U)+(w+U)=(v+w)+U,\qquad
a(v+U)=av+U.
\]
所得线性空间称为 $V$ 关于 $U$ 的\emph{商空间（quotient space）}。
\end{definition}

首先要验证运算不依赖代表元。由 $v+U=w+U$ 当且仅当 $v-w\in U$，若 $v-v'\in U$、$w-w'\in U$，就有 $(v+w)-(v'+w')\in U$ 及 $av-av'\in U$，所以运算良定义。其余线性空间公理由 $V$ 继承，零向量为陪集 $U=0+U$。任意两个陪集或者相等，或者不相交；一个陪集本身是 $V$ 的子空间，当且仅当它就是 $U$。

\begin{theorem}
\label{thm:laex-quotient-complement}
若 $V=U\oplus W$，则 $T:W\to V/U$，$T(w)=w+U$ 是线性同构。特别地，
\[
\dim(V/U)=\dim V-\dim U.
\]
\end{theorem}

\begin{proof}
陪集运算保证 $T$ 线性。若 $T(w)=U$，则 $w\in U\cap W=\{0\}$，故为单射。任意 $v\in V$ 唯一分解为 $u+w$，于是 $v+U=w+U$，故为满射。补空间维数为 $\dim V-\dim U$，从而得到维数公式。若 $e_1,\ldots,e_k$ 是 $U$ 的基，扩张为 $e_1,\ldots,e_n$ 后，$e_{k+1}+U,\ldots,e_n+U$ 就是商空间的一组基。
\end{proof}

商空间把相差一个 $U$ 中向量的元素视为同一个元素，补空间则从每个陪集中选出唯一代表。前者由 $U\subseteq V$ 自然确定，后者依赖选择；因此它们虽然同构，却不是同一个集合。




\section{矩阵的秩}
\label{sec:laex-matrix-rank}

\subsection{基本估计与乘积秩不等式}

秩既可以由主元或非零子式确定，也可以理解为列空间的维数或“未知数个数减去解空间维数”。三种观点对应不同的证明方法：初等变换适合显式分块结构，解空间包含适合比较核，生成空间包含适合比较像。选用哪一种，取决于题目给出的条件最直接地控制了什么。

\begin{proposition}
\label{prop:laex-rank-basic}
在矩阵尺寸相容的条件下，有
\begin{align}
\operatorname{rank}(AB)&\le\min\{\operatorname{rank}A,\operatorname{rank}B\},
\label{eq:laex-product-upper}\\
\operatorname{rank}(A+B)&\le\operatorname{rank}A+\operatorname{rank}B,
\label{eq:laex-sum-upper}\\
\bigl|\operatorname{rank}A-\operatorname{rank}B\bigr|
&\le\operatorname{rank}(A-B).
\label{eq:laex-difference-lower}
\end{align}
此外，可逆矩阵左乘、右乘及非零标量数乘均不改变秩，且
\begin{equation}
\label{eq:laex-block-diagonal-rank}
\operatorname{rank}\begin{pmatrix}A&0\\0&B\end{pmatrix}
=\operatorname{rank}A+\operatorname{rank}B,
\qquad
\operatorname{rank}\begin{pmatrix}A&C\\0&B\end{pmatrix}
\ge\operatorname{rank}A+\operatorname{rank}B.
\end{equation}
\end{proposition}

\begin{proof}
$AB$ 的列是 $A$ 的列的线性组合，故其秩不超过 $\operatorname{rank}A$；把 $B$ 的列用其一个极大无关组表示后，左乘 $A$ 得到 $AB$ 的生成组，故也不超过 $\operatorname{rank}B$。又 $\operatorname{im}(A+B)\subseteq\operatorname{im}A+\operatorname{im}B$，得到和的上界。对 $A=(A-B)+B$ 与 $B=(B-A)+A$ 分别使用该上界，得到差的下界。可逆乘法的结论可由乘积上界在正反两个方向应用得到。

分块对角阵的两组列占据互不重叠的分量，秩相加。对分块上三角阵，分别从 $A,B$ 中选取极大无关列，取它们在整个分块矩阵中的对应列。一个为零的线性组合先看下半部，$B$ 的无关性迫使第二组系数为零，再看上半部，$A$ 的无关性迫使其余系数为零。因此至少有 $\operatorname{rank}A+\operatorname{rank}B$ 个无关列。
\end{proof}

\begin{theorem}
\label{thm:laex-sylvester-frobenius}
若 $A\in M_{m\times n}(\mathbb F)$、$B\in M_{n\times p}(\mathbb F)$，则西尔维斯特（Sylvester）秩不等式为
\begin{equation}
\label{eq:laex-sylvester}
\operatorname{rank}(AB)\ge\operatorname{rank}A+\operatorname{rank}B-n.
\end{equation}
若另有 $C\in M_{p\times q}(\mathbb F)$，则弗罗贝尼乌斯（Frobenius）秩不等式为
\begin{equation}
\label{eq:laex-frobenius}
\operatorname{rank}(ABC)\ge
\operatorname{rank}(AB)+\operatorname{rank}(BC)-\operatorname{rank}B.
\end{equation}
\end{theorem}

\begin{proof}
先证明一个更准确的等式：
\begin{equation}
\label{eq:laex-product-rank-loss}
\operatorname{rank}(AB)
=\operatorname{rank}B-\dim(\operatorname{im}B\cap\ker A).
\end{equation}
取 $\operatorname{im}B\cap\ker A$ 的基并将其扩张为 $\operatorname{im}B$ 的基。交空间的基被 $A$ 送到零，新增基向量的像则无关：否则其某个非零组合落在 $\ker A$ 中，又落在 $\operatorname{im}B$ 中，与扩张后的无关性矛盾。这些像生成 $\operatorname{im}(AB)$，故得到等式。由交空间维数至多为 $\dim\ker A=n-\operatorname{rank}A$，即得 Sylvester 不等式。

分别对 $B$ 和 $BC$ 使用式~\eqref{eq:laex-product-rank-loss}，注意到 $\operatorname{im}(BC)\subseteq\operatorname{im}B$，因此
\[
\begin{aligned}
\operatorname{rank}(AB)-\operatorname{rank}(ABC)
&=\operatorname{rank}B-\operatorname{rank}(BC)\\
&\quad-\dim(\operatorname{im}B\cap\ker A)
+\dim(\operatorname{im}(BC)\cap\ker A)\\
&\le\operatorname{rank}B-\operatorname{rank}(BC).
\end{aligned}
\]
移项就得到 Frobenius 不等式。
\end{proof}

式~\eqref{eq:laex-product-rank-loss} 表明，左乘 $A$ 所损失的维数正是 $B$ 的像中被 $A$ 消去的部分。因此
\begin{equation}
\label{eq:laex-product-rank-equality}
\operatorname{rank}(AB)=\operatorname{rank}B
\quad\Longleftrightarrow\quad
\operatorname{im}B\cap\ker A=\{0\}
\quad\Longleftrightarrow\quad
\ker(AB)=\ker B.
\end{equation}
最后一个等价性也可由 $\ker B\subseteq\ker(AB)$ 与秩与零度公式得到。此时对任何尺寸相容的 $C$，都有 $\operatorname{rank}(ABC)=\operatorname{rank}(BC)$。特别地，$AB=0$ 蕴含 $\operatorname{rank}A+\operatorname{rank}B\le n$。反复应用 Sylvester 不等式还得到，对 $n$ 阶方阵 $A_1,\ldots,A_k$，
\[
\sum_{i=1}^k\operatorname{rank}A_i
\le(k-1)n+\operatorname{rank}(A_1\cdots A_k).
\]

\begin{example}
\label{ex:laex-commuting-rank}
设 $A,B\in M_n(\mathbb F)$ 且 $AB=BA$，证明
\[
\operatorname{rank}(A+B)+\operatorname{rank}(AB)
\le\operatorname{rank}A+\operatorname{rank}B.
\]
\end{example}

\begin{solve}
从列空间看，$\operatorname{im}(A+B)\subseteq\operatorname{im}A+\operatorname{im}B$，而交换条件保证 $\operatorname{im}(AB)\subseteq\operatorname{im}A\cap\operatorname{im}B$。将两个维数上界相加，再用式~\eqref{eq:laex-dimension-sum} 就得到结论。

从解空间看，$\ker A\cap\ker B\subseteq\ker(A+B)$，而 $\ker A+\ker B\subseteq\ker(AB)$，后一包含同样使用 $AB=BA$。将包含关系的维数不等式相加，再用交和维数公式与式~\eqref{eq:laex-rank-nullity}，得到同一结论。两个证明分别刻画列空间与核空间，不能省去交换假设。
\end{solve}

\subsection{分块消元与特殊矩阵}

\begin{proposition}
\label{prop:laex-schur-rank}
设 $M=\begin{pmatrix}A&B\\C&D\end{pmatrix}$，其中 $A$ 为 $r$ 阶可逆方阵，则
\begin{equation}
\label{eq:laex-schur-rank}
\operatorname{rank}M=r+\operatorname{rank}(D-CA^{-1}B).
\end{equation}
矩阵 $D-CA^{-1}B$ 称为关于 $A$ 的舒尔补（Schur complement，亦写作 Schur 补）。若 $D$ 可逆，则也有 $\operatorname{rank}M=\operatorname{rank}D+\operatorname{rank}(A-BD^{-1}C)$。
\end{proposition}

\begin{proof}
以下左右乘矩阵均可逆：
\[
\begin{pmatrix}I&0\\-CA^{-1}&I\end{pmatrix}
\begin{pmatrix}A&B\\C&D\end{pmatrix}
\begin{pmatrix}I&-A^{-1}B\\0&I\end{pmatrix}
=\begin{pmatrix}A&0\\0&D-CA^{-1}B\end{pmatrix}.
\]
再用分块对角阵秩相加公式即可。关于 $D$ 的公式先交换分块行列再同样处理。
\end{proof}

分块初等变换中，乘法顺序由矩阵尺寸决定：左乘系数对应行操作，右乘系数对应列操作。不能像标量一样随意交换各因子。Schur 补公式中的逆也必须由可逆条件保证，不能直接用于奇异分块。

\begin{proposition}
\label{prop:laex-idempotent-involution-rank}
对任意 $A\in M_n(\mathbb F)$，有
\begin{align}
\operatorname{rank}A+\operatorname{rank}(I_n-A)
&=n+\operatorname{rank}(A-A^2),
\label{eq:laex-idempotent-rank}\\
\operatorname{rank}(I_n+A)+\operatorname{rank}(I_n-A)
&=n+\operatorname{rank}(I_n-A^2).
\label{eq:laex-involution-rank}
\end{align}
因此 $A$ 为幂等矩阵（idempotent matrix），即 $A^2=A$，当且仅当第一式左端为 $n$；$A$ 为对合矩阵（involutory matrix），即 $A^2=I_n$，当且仅当第二式左端为 $n$。
\end{proposition}

\begin{proof}
用 $\sim$ 表示相抵变换，有
\[
\begin{pmatrix}A&0\\0&I_n-A\end{pmatrix}
\sim\begin{pmatrix}A&I_n-A\\0&I_n-A\end{pmatrix}
\sim\begin{pmatrix}I_n&I_n-A\\I_n-A&I_n-A\end{pmatrix}
\sim\begin{pmatrix}I_n&0\\0&A-A^2\end{pmatrix}.
\]
最后一步先以第一分块行消去左下块，再消去右上块。比较两侧秩即得第一式。将其中 $A$ 换成 $(I_n+A)/2$，并使用非零数乘不改变秩，得到第二式。
\end{proof}

\subsection{由核空间证明秩与非奇异性}

\begin{proposition}
\label{prop:laex-gram-rank}
若 $A$ 为实矩阵，则
\[
\operatorname{rank}(A^{\mathsf T}A)
=\operatorname{rank}(AA^{\mathsf T})=\operatorname{rank}A.
\]
若 $A$ 为复矩阵，则相应结论应写成 $\operatorname{rank}(A^{\mathsf H}A)=\operatorname{rank}(AA^{\mathsf H})=\operatorname{rank}A$。
\end{proposition}

\begin{proof}
实数情形中，$Ax=0$ 蕴含 $A^{\mathsf T}Ax=0$。反过来，若后者成立，则
\[
0=x^{\mathsf T}A^{\mathsf T}Ax
=(Ax)^{\mathsf T}(Ax)=\sum_i(Ax)_i^2,
\]
每一项非负，所以 $Ax=0$。两矩阵的核相同、列数相同，故秩相等；对 $A^{\mathsf T}$ 使用同一结论得到另一等式。复数情形改用 $x^{\mathsf H}A^{\mathsf H}Ax=\sum_i|(Ax)_i|^2$ 即可。若复数情形仍用转置，结论会失效，例如 $A=(1,\mathrm i)^{\mathsf T}$ 的秩为 $1$，但 $A^{\mathsf T}A=0$。
\end{proof}

\begin{example}
\label{ex:laex-diagonal-dominance}
设 $A=(a_{ij})\in M_n(\mathbb C)$ 满足严格行对角占优条件
\[
|a_{ii}|>\sum_{j\ne i}|a_{ij}|\quad(1\le i\le n).
\]
证明 $A$ 可逆。若 $A$ 为实矩阵且进一步有 $a_{ii}>\sum_{j\ne i}|a_{ij}|$，证明 $\det A>0$。
\end{example}

\begin{solve}
若 $Ax=0$ 有非零解，取 $|x_k|=\max_j|x_j|>0$。第 $k$ 个方程给出
\[
|a_{kk}|\,|x_k|
\le\sum_{j\ne k}|a_{kj}|\,|x_j|
\le\left(\sum_{j\ne k}|a_{kj}|\right)|x_k|,
\]
与严格不等式矛盾，故只有零解，$A$ 可逆。实数正对角元情形令 $D=\operatorname{diag}(a_{11},\ldots,a_{nn})$，考虑 $A(t)=D+t(A-D)$，$0\le t\le1$。路径上始终严格对角占优，行列式连续且处处非零，而 $\det A(0)=\prod_i a_{ii}>0$，由介值定理知其不变号，因此 $\det A>0$。
\end{solve}

\begin{example}
\label{ex:laex-skew-positive}
设 $A$ 为实反对称矩阵，$D$ 为对角元全正的实对角矩阵，证明 $\det(D+A)>0$。
\end{example}

\begin{solve}
若 $(D+tA)x=0$，其中 $t\in\mathbb R$、$x\in\mathbb R^n$，则 $x^{\mathsf T}Ax=0$，故 $0=x^{\mathsf T}Dx=\sum_i d_ix_i^2$，只能 $x=0$。于是实矩阵 $D+tA$ 对所有实数 $t$ 可逆，其行列式是处处非零的连续函数，在 $t=0$ 处为正，因而在 $t=1$ 处也为正。特别地，$\det(I_n\pm A)>0$。
\end{solve}

\subsection{子式、加边子式与伴随矩阵}

\begin{proposition}
\label{prop:laex-bordered-minors}
矩阵 $A$ 的秩为 $r\ge1$，当且仅当存在一个非零的 $r$ 阶子式，且所有 $r+1$ 阶子式都为零。更强地，固定一个非零 $r$ 阶子式后，只需检查包含它的全部 $r+1$ 阶\emph{加边子式（bordered minor）}为零。
\end{proposition}

\begin{proof}
将所选子式移到左上角，写成 $A=\begin{pmatrix}D&B\\C&E\end{pmatrix}$，其中 $D$ 可逆。每个相应加边子式都是
\[
\det\begin{pmatrix}D&b_j\\c_i&e_{ij}\end{pmatrix}
=\det D\,(e_{ij}-c_iD^{-1}b_j).
\]
因此这些加边子式全为零，当且仅当 Schur 补 $E-CD^{-1}B=0$，由式~\eqref{eq:laex-schur-rank} 即得秩为 $r$。必要性来自：任何阶数大于秩的方形子矩阵，其列不可能无关。若秩大于 $r$，从 $r$ 个独立列和相应独立行中可取得更大的非零子式，或直接由上述 Schur 补的非零元素得到非零加边子式。$r=0$ 的情形单独等价于 $A=0$。
\end{proof}

如果分别选取一个秩为 $r$ 的矩阵的 $r$ 个独立行和 $r$ 个独立列，它们交叉处的 $r$ 阶子矩阵必可逆。证明时只需注意所有行都是所选行的线性组合：如果所选列在这些行上的缩短列相关，同一关系就会在全部行上成立，与所选列无关矛盾。对称矩阵或反对称矩阵中，一组独立行的同指标列也独立，因此必有一个阶数等于秩的非零主子式。反对称矩阵的奇数阶行列式满足 $\det D=\det(-D^{\mathsf T})=-\det D$，只能为零，于是\emph{反对称矩阵的秩必为偶数}。

\begin{proposition}
\label{prop:laex-principal-border}
设对称矩阵 $A$ 有一个非零 $r$ 阶主子式 $\det D$，且所有包含它的 $r+1$ 阶及 $r+2$ 阶加边主子式均为零，则 $\operatorname{rank}A=r$。若 $A$ 为反对称矩阵，则只需检查包含该主子式的全部 $r+2$ 阶加边主子式为零。
\end{proposition}

\begin{proof}
将 $D$ 同时置于左上角，并记 Schur 补为 $S$。对称情形下 $S$ 仍对称，$r+1$ 阶加边主子式为零使 $s_{ii}=0$，$r+2$ 阶加边主子式为零再给出
\[
0=\det\begin{pmatrix}0&s_{ij}\\s_{ij}&0\end{pmatrix}=-s_{ij}^2,
\]
故 $S=0$。反对称情形下 $S$ 仍反对称，已自动有 $s_{ii}=0$，相应二阶主子式等于 $s_{ij}^2$，同样迫使 $S=0$。两种情形都由 Schur 补秩公式得到结论；若剩余维数小于二，只需应用存在的相应条件。
\end{proof}

\begin{proposition}
\label{prop:laex-adjugate-rank}
设 $A\in M_n(\mathbb F)$ 且 $n\ge2$，则
\[
\operatorname{rank}\operatorname{adj}(A)=
\begin{cases}
n,&\operatorname{rank}A=n,\\
1,&\operatorname{rank}A=n-1,\\
0,&\operatorname{rank}A\le n-2.
\end{cases}
\]
\end{proposition}

\begin{proof}
若 $A$ 可逆，$\operatorname{adj}(A)=(\det A)A^{-1}$，故满秩。若 $\operatorname{rank}A=n-1$，至少有一个 $n-1$ 阶子式非零，故伴随矩阵非零；又 $A\operatorname{adj}(A)=0$，使其全部列都落在一维核空间中，故秩为 $1$。若 $\operatorname{rank}A\le n-2$，全部 $n-1$ 阶子式为零，伴随矩阵即为零。
\end{proof}

当 $\operatorname{rank}A=n-1$ 时，伴随矩阵的任意非零列就构成 $Ax=0$ 的基础解系。还可用这一点证明：若矩阵每行、每列的元素和都为零，则所有代数余子式相等。秩低于 $n-1$ 时它们全为零；秩为 $n-1$ 时，右核与左核都由 $\mathbf1$ 生成，伴随矩阵的各列与各行都相应成比例，所以它只能是 $c\mathbf1\mathbf1^{\mathsf T}$ 的形式。

\section{相抵标准型及其应用}
\label{sec:laex-equivalence-form}

\subsection{标准型、单侧逆与满秩分解}

\begin{theorem}
\label{thm:laex-equivalence-form}
若 $A\in M_{m\times n}(\mathbb F)$ 的秩为 $r$，则存在可逆矩阵 $P\in M_m(\mathbb F)$、$Q\in M_n(\mathbb F)$，使
\begin{equation}
\label{eq:laex-equivalence-form}
PAQ=J_r=\begin{pmatrix}I_r&0\\0&0\end{pmatrix}.
\end{equation}
同型矩阵相抵，当且仅当秩相同。
\end{theorem}

\begin{proof}
先用行变换得到阶梯形，再用列交换把主元列移到前面，以可逆行列变换把对应可逆分块化为 $I_r$ 并消去其余块，即得到 $J_r$。反过来，左右乘可逆矩阵保持秩，所以同型矩阵有相同标准型当且仅当它们的秩相同。
\end{proof}

相抵允许独立选取左右两个可逆矩阵；相似变换则要求它们互为逆矩阵。因此相抵标准型用于秩问题，不能据此断言矩阵相似于 $J_r$，也不能直接从相抵标准型读出原矩阵的迹或特征值。

\begin{proposition}
\label{prop:laex-one-sided-inverse}
$A\in M_{m\times n}(\mathbb F)$ 存在左逆 $L\in M_{n\times m}(\mathbb F)$、$LA=I_n$，当且仅当 $\operatorname{rank}A=n$；存在右逆 $R\in M_{n\times m}(\mathbb F)$、$AR=I_m$，当且仅当 $\operatorname{rank}A=m$。列满秩矩阵适合左消去律，行满秩矩阵适合右消去律。
\end{proposition}

\begin{proof}
存在单侧逆时由乘积秩上界得到相应满秩性。若列满秩，则标准型为 $PAQ=\binom{I_n}{0}$，取 $L=Q(I_n\mid0)P$ 即有 $LA=I_n$。若行满秩，则 $PAQ=(I_m\mid0)$，取 $R=Q\binom{I_m}{0}P$ 即可。由 $AX=AY$ 左乘 $L$ 得 $X=Y$，右消去律同理。
\end{proof}

\begin{theorem}
\label{thm:laex-rank-factorization}
设 $A\in M_{m\times n}(\mathbb F)$ 的秩为 $r>0$，则存在列满秩的 $B\in M_{m\times r}(\mathbb F)$ 与行满秩的 $C\in M_{r\times n}(\mathbb F)$，使 $A=BC$，称为\emph{满秩分解（rank factorization）}。若还有 $A=B_1C_1=B_2C_2$ 两个满秩分解，则存在 $r$ 阶可逆矩阵 $S$，使
\[
B_2=B_1S,\qquad C_2=S^{-1}C_1.
\]
\end{theorem}

\begin{proof}
从 $A$ 的列中选一个极大无关组组成 $B$，把每一列在该组下的表示系数组成 $C$，便得到 $A=BC$。$C$ 在所选列的对应位置含 $I_r$，所以行满秩。对两个分解，$\operatorname{im}B_1=\operatorname{im}A=\operatorname{im}B_2$，故 $B_1,B_2$ 的列是同一空间的两组基，存在可逆过渡矩阵 $S$ 使 $B_2=B_1S$。代回 $B_1C_1=B_1SC_2$，左消去 $B_1$，得到 $C_2=S^{-1}C_1$。
\end{proof}

满秩分解还表明，秩为 $r$ 的矩阵恰可写成 $r$ 个秩一矩阵之和，且不能用更少个秩一矩阵相加得到。若 $B=(b_1,\ldots,b_r)$，$C$ 的各行为 $c_i^{\mathsf T}$，则 $A=\sum_i b_ic_i^{\mathsf T}$，每项的两因子均非零，秩为 $1$；若少于 $r$ 项，则秩的次可加性会给出矛盾。零矩阵对应空和。

\begin{example}
\label{ex:laex-generalized-inverse}
对任意 $A\in M_{m\times n}(\mathbb F)$，证明存在 $G\in M_{n\times m}(\mathbb F)$，使 $AGA=A$。
\end{example}

\begin{solve}
当 $A=0$ 时取 $G=0$。否则取满秩分解 $A=BC$，由命题~\ref{prop:laex-one-sided-inverse} 选取 $B$ 的左逆 $L$ 与 $C$ 的右逆 $R$，令 $G=RL$，则
\[
AGA=BCRLBC=B(CR)(LB)C=BC=A.
\]
这里依次使用 $CR=I_r$ 与 $LB=I_r$。这个条件称为一种广义逆条件（generalized inverse condition），不要求 $AG$ 或 $GA$ 为单位矩阵，也不保证 $G$ 唯一。
\end{solve}

\subsection{幂等矩阵与矩阵方程}

\begin{proposition}
\label{prop:laex-idempotent-factorization}
设 $A=BC$ 是秩为 $r>0$ 的方阵 $A$ 的满秩分解，则 $A^2=A$ 当且仅当 $CB=I_r$。此时 $\operatorname{tr}A=r$，且 $\mathbb F^n=\operatorname{im}A\oplus\ker A$。
\end{proposition}

\begin{proof}
由 $A^2=A$ 得 $B(CB-I_r)C=0$，用 $B$ 的左逆和 $C$ 的右逆消去两侧即得 $CB=I_r$；反向代入即可。逐元素求和给出 $\operatorname{tr}(BC)=\operatorname{tr}(CB)=r$。任意 $x$ 分解为 $Ax+(x-Ax)$，两项分别属于像和核；若 $v=Au$ 且 $Av=0$，又有 $Av=A^2u=Au=v$，所以交只有零向量。$A=0$ 时迹和直和结论也成立。
\end{proof}

\begin{example}
\label{ex:laex-rectangular-product}
设 $A$ 为 $3\times2$ 矩阵，$B$ 为 $2\times3$ 矩阵，且
\[
AB=C=\begin{pmatrix}8&2&-2\\2&5&4\\-2&4&5\end{pmatrix}.
\]
求 $BA$。
\end{example}

\begin{solve}
直接计算得到 $\operatorname{rank}C=2$ 及 $C^2=9C$。由乘积秩上界与尺寸限制知 $A$ 列满秩、$B$ 行满秩。于是 $A(BA-9I_2)B=0$，用单侧逆消去 $A,B$，得到 $BA=9I_2$。矩形因子的消去依赖满秩条件，不能仅从 $A,B$ 非零推断可消去。
\end{solve}

\begin{proposition}
\label{prop:laex-two-sided-equation}
设 $A\in M_{m\times n}(\mathbb F)$、$B\in M_{p\times q}(\mathbb F)$、$C\in M_{m\times q}(\mathbb F)$。矩阵方程
\[
AX+YB=C,\qquad
X\in M_{n\times q}(\mathbb F),\quad Y\in M_{m\times p}(\mathbb F),
\]
有解，当且仅当
\[
\operatorname{rank}\begin{pmatrix}A&C\\0&B\end{pmatrix}
=\operatorname{rank}A+\operatorname{rank}B.
\]
\end{proposition}

\begin{proof}
若有解，以第一分块列右乘 $-X$ 加到第二分块列，再以第二分块行左乘 $-Y$ 加到第一分块行，即得到分块对角阵，故秩相加。

反过来，将 $A,B$ 分别作相抵变换化为 $J_r,J_s$。若 $A'=PAQ$、$B'=RBS$、$C'=PCS$，则 $X'=Q^{-1}XS$、$Y'=PYR^{-1}$ 使新旧方程等价，相关分块矩阵的秩也不变。因此只需处理标准型。按前 $r$ 行和前 $s$ 列将 $C'$ 分块为 $\begin{pmatrix}C_{11}&C_{12}\\C_{21}&C_{22}\end{pmatrix}$，消去 $J_r,J_s$ 对应的块后，整个矩阵秩为 $r+s+\operatorname{rank}C_{22}$。假设使 $C_{22}=0$，而 $J_rX'$ 可以任取前 $r$ 行，$Y'J_s$ 可以任取前 $s$ 列，故可令前者覆盖 $C_{11},C_{12}$，后者覆盖 $C_{21}$，构造出解。
\end{proof}

\section{线性方程组的解及其应用}
\label{sec:laex-linear-systems}

\subsection{可解性、解的结构与参数讨论}

\begin{theorem}
\label{thm:laex-system-structure}
设 $A\in M_{m\times n}(\mathbb F)$、$b\in\mathbb F^m$、$\operatorname{rank}A=r$。方程组 $Ax=b$ 有解，当且仅当 $\operatorname{rank}(A\mid b)=r$。有解时，若 $x_0$ 是一个特解，$\eta_1,\ldots,\eta_{n-r}$ 是 $\ker A$ 的一组基，则全部解为
\begin{equation}
\label{eq:laex-affine-solution}
x=x_0+\sum_{i=1}^{n-r}t_i\eta_i,\qquad t_i\in\mathbb F.
\end{equation}
其中 $\eta_1,\ldots,\eta_{n-r}$ 称为齐次方程组 $Ax=0$ 的\emph{基础解系（fundamental system of solutions）}。当 $r=n$ 时解唯一，当 $r<n$ 时解有无穷多个。
\end{theorem}

\begin{proof}
可解性等价于 $b$ 属于 $A$ 的列空间，即增添 $b$ 不增加列秩。若 $x$ 与 $x_0$ 都为解，则 $A(x-x_0)=0$；反过来，特解加上任意核向量仍为解。由式~\eqref{eq:laex-rank-nullity}，核空间维数为 $n-r$，因此得到全部解的表达式。唯一性等价于核为零；核非零时，由数域无限可得到无穷多个解。
\end{proof}

非齐次解集 $x_0+\ker A$ 是一个\emph{仿射子空间（affine subspace）}，当 $b\ne0$ 时不含零向量，因而不是线性子空间。若 $x^{(1)},\ldots,x^{(k)}$ 都是解，则
\[
A\left(\sum_i c_ix^{(i)}\right)=\left(\sum_i c_i\right)b.
\]
当 $b\ne0$ 时，该组合仍为解当且仅当 $\sum_i c_i=1$；齐次情形则对系数没有这一限制。基础解系也不是任取 $n-r$ 个解向量，而必须是 $n-r$ 个线性无关的齐次解。两个基础解系互相通过一个可逆的 $n-r$ 阶系数矩阵联系，本质上就是同一核空间的基变换。

带参数的消元尤其要区分\emph[caution]{可能为零的主元}。若某一步拟除以 $g(\lambda)$，必须先分成 $g(\lambda)=0$ 与 $g(\lambda)\ne0$ 两种情况，零点处回到原矩阵重新计算；不能把一般参数情形的阶梯形直接代入所有特殊值。另一种办法是利用结构找出参数导致秩变化的位置，再逐一比较系数矩阵和增广矩阵的秩。

\begin{example}
\label{ex:laex-parameter-consistency}
讨论含参数 $\lambda\in\mathbb F$ 的方程组
\[
\begin{cases}
2x_1+3x_2+x_3+x_4=1,\\
x_1+2x_2-x_3+4x_4=2,\\
x_1+3x_2-4x_3+11x_4=\lambda.
\end{cases}
\]
\end{example}

\begin{solve}
第三个系数行等于第二行的三倍减去第一行，因此有解必须且只需 $\lambda=3\cdot2-1=5$。前两行独立，系数矩阵秩为 $2$。当 $\lambda=5$ 时令 $x_3=s$、$x_4=t$，得到
\[
x=\begin{pmatrix}-4\\3\\0\\0\end{pmatrix}
+s\begin{pmatrix}-5\\3\\1\\0\end{pmatrix}
+t\begin{pmatrix}10\\-7\\0\\1\end{pmatrix},\qquad s,t\in\mathbb F.
\]
当 $\lambda\ne5$ 时，增广矩阵秩为 $3$ 而系数矩阵秩为 $2$，所以无解。
\end{solve}

\begin{example}
\label{ex:laex-parameter-symmetric-system}
讨论参数 $k\in\mathbb F$ 下的方程组
\[
\begin{cases}
kx_1+x_2+x_3=-2,\\
x_1+kx_2+x_3=-2,\\
x_1+x_2+kx_3=-2.
\end{cases}
\]
\end{example}

\begin{solve}
任意两式相减得到 $(k-1)(x_i-x_j)=0$，三式相加得到 $(k+2)(x_1+x_2+x_3)=-6$。若 $k\ne1,-2$，则 $x_1=x_2=x_3=-2/(k+2)$，为唯一解。若 $k=1$，三个方程相同，通解为
\[
x=\begin{pmatrix}-2\\0\\0\end{pmatrix}
+s\begin{pmatrix}-1\\1\\0\end{pmatrix}
+t\begin{pmatrix}-1\\0\\1\end{pmatrix},\qquad s,t\in\mathbb F.
\]
若 $k=-2$，三式相加得到 $0=-6$，故无解。这里两个特殊值分别对应失去差分约束与出现不相容的总和约束。
\end{solve}

\subsection{解空间、行空间与可解性的另一种判据}

\begin{proposition}
\label{prop:laex-kernel-rowspace}
设 $A\in M_{m\times n}(\mathbb F)$，$c\in\mathbb F^n$。若 $Ax=0$ 的每个解都满足 $c^{\mathsf T}x=0$，则 $c^{\mathsf T}$ 是 $A$ 的行向量的线性组合；反之亦然。因此，两个含有相同未知数的齐次方程组同解，当且仅当其系数矩阵的行空间相同。
\end{proposition}

\begin{proof}
向 $A$ 添加行 $c^{\mathsf T}$ 得到矩阵 $D$。题设等价于 $\ker D=\ker A$，由秩与零度公式得到 $\operatorname{rank}D=\operatorname{rank}A$。添加的一行不能增加行空间维数，所以它已在原行空间中。反向结论只需对已有各方程取相应线性组合。同解的两个方程组逐行使用上述结论，即得到行空间互相包含。
\end{proof}

\begin{theorem}
\label{thm:laex-solvability-alternative}
设 $A\in M_{m\times n}(\mathbb F)$、$b\in\mathbb F^m$。方程组 $Ax=b$ 有解，当且仅当
\begin{equation}
\label{eq:laex-solvability-dual}
A^{\mathsf T}y=0\quad\Longrightarrow\quad b^{\mathsf T}y=0.
\end{equation}
等价地，下面两个方程组恰有一个可解：
\[
Ax=b;\qquad
A^{\mathsf T}y=0,\quad b^{\mathsf T}y=1.
\]
\end{theorem}

\begin{proof}
若 $Ax=b$，则对任意 $A^{\mathsf T}y=0$ 有 $b^{\mathsf T}y=x^{\mathsf T}A^{\mathsf T}y=0$。反过来，将命题~\ref{prop:laex-kernel-rowspace} 用于矩阵 $A^{\mathsf T}$，题设就使 $b^{\mathsf T}$ 属于其行空间，即 $b\in\operatorname{im}A$，故有解。若 $Ax=b$ 无解，上述蕴含失效，存在 $A^{\mathsf T}y=0$ 且 $b^{\mathsf T}y\ne0$，将 $y$ 除以这个非零标量得到第二组的解；两组同时有解则会推出 $0=1$。
\end{proof}

这个判据使用的是代数配对 $b^{\mathsf T}y$，在任意数域上都成立，不依赖正定性。在实数域可解释为 $b$ 与左核中的每个向量正交；在复数域若采用通常的 Hermite 内积，则应相应将条件写成 $A^{\mathsf H}y=0\Rightarrow y^{\mathsf H}b=0$，不能把两种记号随意混合。

任意子空间 $U\subseteq\mathbb F^n$ 都可以表示为某个齐次方程组的解空间。若基矩阵 $B\in M_{n\times r}(\mathbb F)$ 满足 $U=\operatorname{im}B$，先求 $B^{\mathsf T}y=0$ 的基础解系，并把其转置排成矩阵 $A$ 的各行，则 $AB=0$、$\operatorname{rank}A=n-r$。于是 $U\subseteq\ker A$ 且两空间维数同为 $r$，所以 $U=\ker A$。$r=n$ 时可取 $A$ 为零矩阵；$r=0$ 时可取 $A=I_n$。这一构造把“生成式”与“方程式”两种子空间描述联系起来。

\subsection{矩阵方程与公共解}

\begin{proposition}
\label{prop:laex-matrix-system}
给定 $A\in M_{m\times n}(\mathbb F)$、$B\in M_{m\times p}(\mathbb F)$，矩阵方程 $AX=B$ 有解，当且仅当 $\operatorname{rank}(A\mid B)=\operatorname{rank}A$。若 $X_0$ 是一个解，$N$ 的列组成 $\ker A$ 的一组基，则全部解为
\[
X=X_0+NZ,\qquad Z\in M_{(n-r)\times p}(\mathbb F),\quad r=\operatorname{rank}A.
\]
\end{proposition}

\begin{proof}
按列写成 $Ax_j=b_j$，可解性等价于 $B$ 的每一列都落在 $A$ 的列空间中。每一列的齐次部分可唯一写成 $Nz_j$，将这些系数列拼成 $Z$ 即得通解。若 $r=n$，齐次部分为空，解唯一。
\end{proof}

例如，对尺寸相容的 $A,B$，存在 $C$ 使 $ABC=A$，当且仅当 $\operatorname{rank}(AB)=\operatorname{rank}A$。必要性由乘积秩上界得到；充分性由 $\operatorname{im}(AB)\subseteq\operatorname{im}A$ 且维数相同得到两像空间相等，再逐列求解 $(AB)C=A$。

齐次方程组 $Ax=0$、$Bx=0$ 的公共解空间为
\[
\ker A\cap\ker B=\ker\binom AB,
\]
所以公共解空间维数为 $n-\operatorname{rank}\binom AB$。若只给出两个非齐次方程组的通解
\[
x=p+Nu,\qquad x=q+Mv,
\]
则求公共解等价于对参数求解
\begin{equation}
\label{eq:laex-common-affine}
(N\mid-M)\binom uv=q-p.
\end{equation}
解出后代回任一通解即可。若存在一个公共解 $x_*$，全部公共解构成 $x_*+(\operatorname{im}N\cap\operatorname{im}M)$；否则交集为空。两个非齐次解集即使各自非空，也可能没有公共解，因此必须先检查参数方程的相容性。

\subsection{解析几何中的秩判定}

设 $p_1,\ldots,p_s\in\mathbb R^d$，其仿射包是
\[
p_1+\operatorname{span}(p_2-p_1,\ldots,p_s-p_1),
\]
所以仿射维数就是差向量矩阵的秩。特别地，平面点共线当且仅当差向量的秩至多为 $1$，空间点共面当且仅当差向量的秩至多为 $2$。这里使用“至多”，允许点重合或落在更低维的集合中。若把点坐标作为列并在底部补一行 $1$，则其矩阵秩等于仿射维数加 $1$：将其余列减去第一列后，差向量的末分量为零，而第一列末分量为 $1$，恰多出一维。

\begin{example}
\label{ex:laex-circle}
设平面上三点 $(x_i,y_i)$，$1\le i\le3$，不共线，求经过它们的圆的方程，并给出第四点与它们共圆的判据。
\end{example}

\begin{solve}
圆可设为 $x^2+y^2+ux+vy+w=0$，代入三点得到关于 $u,v,w$ 的线性方程组。不共线保证系数矩阵 $(x_i,y_i,1)_{i=1}^3$ 可逆，故唯一确定圆，等价的行列式形式为
\begin{equation}
\label{eq:laex-circle-determinant}
\det\begin{pmatrix}
x^2+y^2&x&y&1\\
x_1^2+y_1^2&x_1&y_1&1\\
x_2^2+y_2^2&x_2&y_2&1\\
x_3^2+y_3^2&x_3&y_3&1
\end{pmatrix}=0.
\end{equation}
第一项的系数不为零，所以不会退化为直线；配方后半径也必为正，否则不能经过三个不共线实点。第四点共圆当且仅当将其坐标代入式~\eqref{eq:laex-circle-determinant} 后行列式为零。若原来的点全共线，这个行列式条件可能描述退化情形，不能无条件当作圆的判据。
\end{solve}

\begin{example}
\label{ex:laex-conic-pencil}
设平面四个不同点中任意三点不共线，$q_1,q_2$ 是经过它们的两个线性无关的次数至多为二的多项式。证明任意经过这四点的次数至多为二的多项式都可写为 $\lambda q_1+\mu q_2$。
\end{example}

\begin{solve}
全体次数至多为二的二元多项式以 $x^2,xy,y^2,x,y,1$ 为基，维数为 $6$。四点取值给出四个线性条件。为证明这些条件独立，对每个点 $p_i$，选另外三点中的两点确定一条直线，再取经过剩余一点且不经过 $p_i$ 的另一条直线。由于任意三点不共线，两条直线方程的乘积在另外三点为零而在 $p_i$ 非零。适当缩放后，四点的取值为第 $i$ 个标准单位向量，故取值映射的秩为 $4$，核维数为 $2$。$q_1,q_2$ 是此核中的两个无关向量，因而构成基，得到结论。

几何上，$(\lambda,\mu)\ne(0,0)$ 给出过四点的二次曲线束，其中可能含退化曲线。上述点位条件是保证本论证成立的充分条件；若四点全部共线，二次多项式在该直线上的限制只有三个自由系数，四个取值条件便不独立。仅说“两条不同二次曲线有四个公共点”也不足以排除共享直线分支，故应明确这里的点位与无关性条件。
\end{solve}

\section{基础训练：综合方法与易错辨析}
\label{sec:laex-training}

本节将原书基础训练中的代表性问题按方法整合。计算题首先检查秩与尺寸是否匹配，证明题首先寻找可以控制的生成空间、核空间或可逆分块；选择题中的错误选项则常由一个低维反例排除。例如行列式为零只保证行向量组相关，并不保证有零行或有两行成比例；方程比未知数多也并不必然无解；非齐次方程组行满秩则对每个常数向量都有解，但是否唯一还取决于列数。

\begin{example}
\label{ex:laex-independent-functions}
设 $a_1,\ldots,a_n$ 是互异实数，证明 $e^{a_1x},\ldots,e^{a_nx}$ 在实函数空间中线性无关，并说明三角函数组 $1,\sin x,\cos x,\ldots,\sin nx,\cos nx$ 为何也无关。
\end{example}

\begin{solve}
若 $\sum_j c_je^{a_jx}=0$ 对所有实数 $x$ 成立，分别求 $k=0,\ldots,n-1$ 阶导数并令 $x=0$，得到 $\sum_j a_j^kc_j=0$。系数矩阵为非奇异范德蒙德矩阵，因此所有 $c_j=0$。这里只对有限和求导，没有交换无穷过程。

对三角函数，若它们的线性组合恒为零，依次与 $1,\sin kx,\cos kx$ 相乘并在 $[-\pi,\pi]$ 上积分。由积化和差公式可得不同函数乘积的积分为零，而 $1$ 的平方积分为 $2\pi$，其余各函数的平方积分为 $\pi$，于是各系数也逐个为零。所用正交关系在此区间上成立，不能只凭不同函数的外观认定无关。
\end{solve}

\begin{example}
\label{ex:laex-outer-products}
设 $a_1,\ldots,a_p\in\mathbb F^m$ 无关，$b_1,\ldots,b_q\in\mathbb F^n$ 无关，证明 $pq$ 个矩阵 $a_ib_j^{\mathsf T}$ 线性无关。
\end{example}

\begin{solve}
令 $A=(a_1,\ldots,a_p)$、$B=(b_1,\ldots,b_q)$。若 $\sum_{i,j}c_{ij}a_ib_j^{\mathsf T}=0$，写成矩阵即 $ACB^{\mathsf T}=0$，其中 $C=(c_{ij})$。$A$ 列满秩，$B^{\mathsf T}$ 行满秩，故可用单侧逆先后消去它们，得到 $C=0$，所以全部矩阵无关。若缺少任一原向量组的无关性，就不能实施相应消去。
\end{solve}

\begin{example}
\label{ex:laex-biquadratic-field}
求 $\mathbb Q(\sqrt2,\sqrt3)=\{a+b\sqrt2+c\sqrt3+d\sqrt6:a,b,c,d\in\mathbb Q\}$ 在 $\mathbb Q$ 上的维数。
\end{example}

\begin{solve}
先令 $K=\mathbb Q(\sqrt2)$。它以 $1,\sqrt2$ 为有理基，且 $\sqrt3\notin K$：若 $\sqrt3=u+v\sqrt2$，平方并比较有理部分与 $\sqrt2$ 部分得到 $uv=0$；若 $v=0$，则 $u^2=3$，若 $u=0$，则 $v^2=3/2$，两者都不可能由有理数满足，依据是有理数平方中每个素因子的指数均为偶数。因此 $1,\sqrt3$ 在 $K$ 上无关。集合 $K+K\sqrt3$ 对乘法封闭，而且非零元的逆为
\[
(u+v\sqrt3)^{-1}=\frac{u-v\sqrt3}{u^2-3v^2}\quad(u,v\in K),
\]
分母非零，否则会推出 $\sqrt3\in K$ 或 $u=v=0$。它是 $K$ 上维数为 $2$ 的数域，由命题~\ref{prop:laex-scalar-restriction} 得到有理维数为 $4$，基为 $1,\sqrt2,\sqrt3,\sqrt6$。
\end{solve}

\begin{example}
\label{ex:laex-telescoping-rank}
证明 $n$ 阶方阵满足
\[
\operatorname{rank}(AB-I_n)
\le\operatorname{rank}(A-I_n)+\operatorname{rank}(B-I_n).
\]
若 $B_1,\ldots,B_k$ 都为 $n$ 阶幂等矩阵，且 $A=B_1\cdots B_k$，进一步证明
\[
\operatorname{rank}(I_n-A)\le k\bigl(n-\operatorname{rank}A\bigr).
\]
\end{example}

\begin{solve}
第一式由 $AB-I_n=(A-I_n)B+(B-I_n)$ 及秩的乘积上界、和的上界直接得到。第二式使用有限乘积的望远镜分解
\[
I_n-A=\sum_{j=1}^k B_1\cdots B_{j-1}(I_n-B_j),
\]
其中 $j=1$ 时前面的空乘积为 $I_n$。因为 $B_j$ 幂等，式~\eqref{eq:laex-idempotent-rank} 给出 $\operatorname{rank}(I_n-B_j)=n-\operatorname{rank}B_j$，而乘积秩上界给出 $\operatorname{rank}B_j\ge\operatorname{rank}A$。于是各项秩至多为 $n-\operatorname{rank}A$，相加得到所需估计。
\end{solve}

\begin{example}
\label{ex:laex-sum-idempotents}
设 $P=A+B$ 为幂等矩阵，且 $\operatorname{rank}P=\operatorname{rank}A+\operatorname{rank}B$，证明 $A^2=A$、$B^2=B$、$AB=BA=0$。
\end{example}

\begin{solve}
由 $\operatorname{im}P\subseteq\operatorname{im}A+\operatorname{im}B$ 与维数上界知
\[
\operatorname{im}P=\operatorname{im}A\oplus\operatorname{im}B.
\]
幂等矩阵 $P$ 在其像上作用为恒等，所以 $PA=A$、$PB=B$。由前式得到 $(A^2-A)x=-BAx$，左端属于 $\operatorname{im}A$，右端属于 $\operatorname{im}B$，两像交为零，故 $A^2=A$、$BA=0$；对后式同理得到 $B^2=B$、$AB=0$。仅有 $A+B$ 幂等而没有秩相加条件时，不能推出这些结论。
\end{solve}

\begin{example}
\label{ex:laex-normal-rank-stability}
设实方阵 $A$ 满足 $AA^{\mathsf T}=a^2I_n$，其中 $a\in\mathbb R$，证明
\[
\operatorname{rank}\bigl((aI_n-A)^2\bigr)=\operatorname{rank}(aI_n-A).
\]
\end{example}

\begin{solve}
若 $a=0$，由命题~\ref{prop:laex-gram-rank} 知 $A=0$，结论成立。若 $a\ne0$，则 $A$ 可逆，从 $AA^{\mathsf T}=a^2I_n$ 得到 $A^{\mathsf T}=a^2A^{-1}$，因而 $A^{\mathsf T}A=a^2I_n$。令 $M=aI_n-A$，可得 $MM^{\mathsf T}=M^{\mathsf T}M$。以 $\|x\|^2=x^{\mathsf T}x$ 表示实向量的欧几里得范数的平方，则 $\|Mx\|^2=\|M^{\mathsf T}x\|^2$，故 $\ker M=\ker M^{\mathsf T}$。若 $M^2x=0$，令 $v=Mx$，则 $v\in\ker M=\ker M^{\mathsf T}$，从而
\[
\|v\|^2=(Mx)^{\mathsf T}v=x^{\mathsf T}M^{\mathsf T}v=0.
\]
因此 $Mx=0$，证明 $\ker M^2=\ker M$，再由秩与零度公式得到结论。
\end{solve}

\begin{example}
\label{ex:laex-zero-coordinate}
设 $Ax=b$ 有解，记增广矩阵 $\widetilde A=(A\mid b)$，从中删除第 $k$ 个系数列得到 $D$。证明每个解的第 $k$ 个分量都为零，当且仅当 $\operatorname{rank}D<\operatorname{rank}A$。
\end{example}

\begin{solve}
记 $A=(a_1,\ldots,a_n)$。若所有解都满足 $x_k=0$，则 $b$ 可由其余系数列生成；而 $a_k$ 不能由这些列生成，否则将该齐次关系加到一个已有特解上，就得到 $x_k\ne0$ 的解。因此添加 $a_k$ 恰增加一维，$\operatorname{rank}D=\operatorname{rank}A-1$。

反过来，若有一个解满足 $x_k\ne0$，从 $\sum_jx_ja_j=b$ 可解出 $a_k$ 为其余列与 $b$ 的线性组合，所以删去 $a_k$ 不改变增广矩阵的秩。因方程组有解，$\operatorname{rank}\widetilde A=\operatorname{rank}A$，这就与 $\operatorname{rank}D<\operatorname{rank}A$ 矛盾。这里删的是增广矩阵中的系数列，并保留常数列；若仅考察 $A$ 删列，不能反映特解的分量限制。
\end{solve}

\begin{example}
\label{ex:laex-three-lines}
设实平面上三条不同的直线为
\[
ax+by+c=0,\qquad bx+cy+a=0,\qquad cx+ay+b=0.
\]
证明它们交于同一点，当且仅当 $a+b+c=0$。
\end{example}

\begin{solve}
将三式看成关于 $x,y$ 的方程组，其系数矩阵为 $K=\begin{pmatrix}a&b\\b&c\\c&a\end{pmatrix}$。若其秩小于 $2$，所有二阶子式为零，得到 $a^2=bc$、$b^2=ac$、$c^2=ab$，故
\[
(a-b)^2+(b-c)^2+(c-a)^2=0.
\]
实数条件迫使 $a=b=c$，这会使三式相同或不能定义直线，与假设矛盾。因此 $\operatorname{rank}K=2$，三条线有公共点当且仅当增广矩阵秩为 $2$，也即
\[
\det\begin{pmatrix}a&b&c\\b&c&a\\c&a&b\end{pmatrix}
=-(a+b+c)(a^2+b^2+c^2-ab-bc-ca)=0.
\]
第二因子是三个差的平方和的一半，在三条线不同的条件下严格为正，所以条件恰为 $a+b+c=0$。系数秩为 $2$ 还保证公共点唯一。
\end{solve}





